% Generated from locale/tr/content/propositional-logic/syntax-and-semantics/soundness.tex; do not hand-edit.


\olfileid[tr]{pl}{prp}{snd}

\olsection{Önerme Mantığının Sağlamlığı}

\begin{thm}[Sağlamlık]
\ollabel{thm:soundness}
$\Gamma\Proves!A$ ise $\Gamma \Entails !A$ olur.
\end{thm}

\begin{proof}
Teoremler üzerinde tümevarımla. $!A$ bir aksiyomsa $\Entails
!A$ ve dolayısıyla $\Gamma \Entails!A$ olur. Benzer biçimde
$!A \in \Gamma$ ise $\Gamma \Entails!A$ olur. Tümevarım adımı için $!A$'nın,
$!C$ ve $!C\lif!A$'dan \emph{modus ponens} ile elde edildiğini varsayın;
ayrıca tümevarım varsayımı olarak teoremin $!C\lif !A$ ve $!C$ için geçerli
olduğunu kabul edin. \olref[pl][syn][sem]{prop:semanticalfacts} önermesinin üçüncü
maddesi istenen sonucu verir.
\end{proof}

\begin{cor}
$\Gamma$ sağlanabilirse tutarlıdır. Dolayısıyla önerme mantığı tutarlıdır.
\end{cor}

